\pgfdeclarelayer{functionsgrid}
\pgfsetlayers{functionsgrid,main}
\begin{tikzpicture}[x=1.1cm,y=0.56cm,>=Latex,line cap=round]
  \draw[->,black!55] (-4,0)--(5.2,0) node[right] {$x$};
  \draw[->,black!55] (0,-4)--(0,7.2) node[above] {$y$};
  \draw[black!50] (-4,-0.055)--(-4,0.055);
\node[above=4pt,inner sep=1pt] (xt0) at (-4,0) {$-4$};
\draw[black!50] (-3,-0.055)--(-3,0.055);
\node[above=4pt,inner sep=1pt] (xt1) at (-3,0) {$-3$};
\draw[black!50] (-2,-0.055)--(-2,0.055);
\node[above=4pt,inner sep=1pt] (xt2) at (-2,0) {$-2$};
\draw[black!50] (-1,-0.055)--(-1,0.055);
\node[above left=5pt,inner sep=1pt] (xt3) at (-1,0) {$-1$};
\draw[black!50] (1,-0.055)--(1,0.055);
\node[below=4pt,inner sep=1pt] (xt4) at (1,0) {$1$};
\draw[black!50] (2,-0.055)--(2,0.055);
\node[below=4pt,inner sep=1pt] (xt5) at (2,0) {$2$};
\draw[black!50] (3,-0.055)--(3,0.055);
\node[below=4pt,inner sep=1pt] (xt6) at (3,0) {$3$};
\draw[black!50] (4,-0.055)--(4,0.055);
\node[below=4pt,inner sep=1pt] (xt7) at (4,0) {$4$};
\draw[black!50] (5,-0.055)--(5,0.055);
\node[below=4pt,inner sep=1pt] (xt8) at (5,0) {$5$};
\draw[black!50] (-0.055,-4)--(0.055,-4);
\node[left=4pt,inner sep=1pt] (yt9) at (0,-4) {$-4$};
\draw[black!50] (-0.055,-2)--(0.055,-2);
\node[left=4pt,inner sep=1pt] (yt11) at (0,-2) {$-2$};
\draw[black!50] (-0.055,2)--(0.055,2);
\node[left=4pt,inner sep=1pt] (yt14) at (0,2) {$2$};
\draw[black!50] (-0.055,4)--(0.055,4);
\node[left=4pt,inner sep=1pt] (yt16) at (0,4) {$4$};
\draw[black!50] (-0.055,6)--(0.055,6);
\node[left=4pt,inner sep=1pt] (yt18) at (0,6) {$6$};

  \begin{scope}
    \clip (-4,-4) rectangle (5,7);
    \draw[black!50,dashed,line width=0.8pt] (-4,-3)--(5,6);
\draw[green!55!black,dashed,line width=0.8pt] (1,-4)--(1,-0.9);
\draw[green!55!black,dashed,line width=0.8pt] (1,0.2)--(1,7);
\draw[THMBlue,line width=1.15pt,domain=-4:0.65,samples=181] plot (\x,{\x+1+2/(\x-1)});
\draw[THMBlue,line width=1.15pt,domain=1.43:5,samples=161] plot (\x,{\x+1+2/(\x-1)});
  \end{scope}
  \node[anchor=south] at (0.500,8.539) {$\displaystyle f(x)=\frac{x^2+1}{x-1}$};
  \node[anchor=north] at (0.5,-5.05) {\Lang{Gestrichelt: die beiden Asymptoten}{Dashed: the two asymptotes}};
  \begin{pgfonlayer}{functionsgrid}
    \begin{scope}[even odd rule]
      \clip (-4,-4) rectangle (5,7)
        (xt0.south west) rectangle (xt0.north east)
        (xt1.south west) rectangle (xt1.north east)
        (xt2.south west) rectangle (xt2.north east)
        (xt3.south west) rectangle (xt3.north east)
        (xt4.south west) rectangle (xt4.north east)
        (xt5.south west) rectangle (xt5.north east)
        (xt6.south west) rectangle (xt6.north east)
        (xt7.south west) rectangle (xt7.north east)
        (xt8.south west) rectangle (xt8.north east)
        (yt9.south west) rectangle (yt9.north east)
        (yt11.south west) rectangle (yt11.north east)
        (yt14.south west) rectangle (yt14.north east)
        (yt16.south west) rectangle (yt16.north east)
        (yt18.south west) rectangle (yt18.north east)
;
      \draw[step=1,black!8,line width=0.25pt] (-4,-4) grid (5,7);
    \end{scope}
  \end{pgfonlayer}
\end{tikzpicture}
